\section{Threat Model}
\label{sec:threat-model}

\tesserahash{} is a \emph{public artifact}: assume an attacker sees every byte of every identifier, can replay them, and may submit arbitrary content as input to a \tesserahash{} computation. The identifier itself carries no secret. The security goals reduce to three properties of the wire format and the SHAKE-256 \citep{fips202} tail.

\subsection{What \tesserahash{} provides}

\begin{enumerate}
  \item \textbf{Collision resistance of the tail.} For any two distinct contents $c_1 \ne c_2$, the probability that $\mathrm{SHAKE\text{-}256}(c_1, 64) = \mathrm{SHAKE\text{-}256}(c_2, 64)$ is bounded by the SHAKE-256 collision-resistance assumption ($2^{256}$ classical, $\sim 2^{170}$ Brassard--H{\o}yer--Tapp \citep{bht1998} quantum).
  \item \textbf{Second-preimage resistance of the tail.} Given a target $h$, the cost of finding any $c$ with $\mathrm{SHAKE\text{-}256}(c, 64) = h$ is $\min(2^{d}, 2^{256}) = 2^{256}$ classical per FIPS 202 \citep{fips202} (the SHAKE-256 sponge capacity caps generic preimage-class security at 256 bits regardless of output length $d$), and $2^{256}$ under Grover \citep{grover1996} on the 512-bit output --- generations beyond any plausible attacker budget.
  \item \textbf{Malleability resistance of the wire encoding.} The 73-byte binary and its canonical Base64URL / Crockford Base32 encodings round-trip byte-identical; canonical decoders reject non-canonical inputs (wrong length, invalid variant byte, padding-bit anomalies in Base32).
\end{enumerate}

\subsection{What \tesserahash{} does not provide}

\begin{itemize}
  \item \textbf{Authentication.} Anyone with the content can recompute the same \tesserahash{}-tail. \tesserahash{} is keyless. If the question is ``did this specific party produce this hash,'' use a signature (PASETO \citep{paseto-v4}, ML-DSA \citep{fips204}, classical Ed25519 \citep{rfc8032}) over the \tesserahash{} or its referenced content. \tesseraid{}'s security analysis makes the same point for identifier authentication; both sibling projects defer producer authentication to an outer envelope.
  \item \textbf{Confidentiality.} The 64-byte SHAKE-256 tail leaks short-content guessability: for low-entropy content, an attacker who can guess the input can verify it by recomputing the tail. This is the same property that breaks ``hash the password and ship the hash'' patterns. If the content is low-entropy and must remain secret, wrap it in \tesseraseal{} (\cref{sec:monorepo}); do not rely on \tesserahash{} for secrecy.
  \item \textbf{Tamper-evidence of the identifier itself.} \tesserahash{} has no MAC or checksum over the full 73-byte form. A bit-flipped \tesserahash{} may decode to another valid \tesserahash{} (with a different tail). If integrity matters, the \tesserahash{} belongs inside a signed envelope or carried as the AEAD AD of a \tesseraseal{} envelope.
  \item \textbf{Producer non-repudiation.} Two parties who agree on the content trivially agree on the tail. Provenance comes from the timestamp/counter \emph{and} the outer protocol, not from the \tesserahash{} alone.
  \item \textbf{Side-channel resistance.} \tesserahash{} is keyless, so timing attacks on the generator are uninteresting (any timing signal an attacker recovers reveals only the content, which they may already control). When \tesserahash{} is used \emph{inside} a security boundary (e.g., as part of a keyed protocol), the surrounding code's side-channel posture is the relevant one.
\end{itemize}

\subsection{Quantum-relevant attacks}

\begin{itemize}
  \item \textbf{Grover's algorithm \citep{grover1996}} reduces generic search cost from $2^n$ to $2^{n/2}$. For \tesserahash{}'s 512-bit tail, that is $2^{256}$ --- coinciding with the FIPS 202 classical second-preimage cap of $\min(2^{d}, 2^{256})$, and astronomical either way.
  \item \textbf{Brassard--H{\o}yer--Tapp \citep{bht1998}} reduces collision cost from $2^{n/2}$ to $2^{n/3}$. For 512 bits, that is $\sim 2^{170}$ --- comfortable but the dominant quantum bound for \tesserahash{}.
  \item \textbf{KyberSlash, PQShield power-analysis, Hertzbleed, GoFetch.} Not applicable: \tesserahash{} has no asymmetric crypto and no secret-dependent control flow. SHAKE-256/Keccak is a fixed-structure permutation \citep{fips202} with no input-dependent timing in stock implementations.
\end{itemize}

\subsection{Random source quality}

\tesserahash{}'s tail is deterministic in the content; no CSPRNG is invoked for the tail. The 16-bit per-millisecond counter advances deterministically by 1 within a millisecond, with no randomness involved. The timestamp comes from the OS wall clock (Unix real-time clock, milliseconds since the epoch); the generator clamps on backward clock jumps so emitted timestamps never regress (\cref{sec:wire-format}). TesseraID's discipline on \texttt{getrandom(2)} / \texttt{getentropy(3)} / \texttt{BCryptGenRandom} does not apply to \tesserahash{} because the format has no random field.
